\subsection*{Operator Precedence and Associativity}

When an expression combines several operators, \emph{precedence} decides
which operator binds first, and \emph{associativity} decides the order in
which two operators of the same precedence group are applied. The table
below lists the operators, function application, the conditional
expression (\textbf{\texttt{if}} / \textbf{\texttt{else}}), and
\textbf{\texttt{lambda}} from the grammar above, in decreasing order of
precedence: an entry binds more tightly than every entry below it. Use
parentheses to override the default grouping.

\begin{center}
  \begin{tabular}{l|l|c}
    Operators & Description & Associativity \\ \hline
    \texttt{\textbf{( )}} & function application & left to right \\
    \texttt{\textbf{**}} & exponentiation & right to left \\
    \texttt{\textbf{+}}\ \texttt{\textbf{-}} & unary plus, unary minus & n/a \\
    \texttt{\textbf{*}}\ \texttt{\textbf{/}}\ \texttt{\textbf{\%}} & multiplication, division, modulo & left to right \\
    \texttt{\textbf{+}}\ \texttt{\textbf{-}} & addition, subtraction & left to right \\
    \texttt{\textbf{<}}\ \texttt{\textbf{>}}\ \texttt{\textbf{==}}\ \texttt{\textbf{!=}}\ \texttt{\textbf{>=}}\ \texttt{\textbf{<=}} & comparison & left to right \\
    \textbf{\texttt{not}} & logical negation & n/a \\
    \textbf{\texttt{and}} & logical conjunction & left to right \\
    \textbf{\texttt{or}} & logical disjunction & left to right \\
    \textbf{\texttt{if}} / \textbf{\texttt{else}} & conditional expression & right to left \\
    \textbf{\texttt{lambda}} & lambda expression & n/a \\
  \end{tabular}
\end{center}

Notes:
\begin{itemize}
\item Since \texttt{\textbf{**}} binds more tightly than unary
  \texttt{\textbf{+}}/\texttt{\textbf{-}}, the expression
  \texttt{\textbf{-2 ** 2}} is \texttt{\textbf{-(2 ** 2)}}, not
  \texttt{\textbf{(-2) ** 2}}.
\item Comparison operators do not chain: \texttt{\textbf{a < b < c}} is
  parsed as \texttt{\textbf{(a < b) < c}}, not as a combined test of
  \texttt{\textbf{a < b}} and \texttt{\textbf{b < c}} as in full Python.
  Since \texttt{\textbf{a < b}} evaluates to a \textbf{\texttt{bool}},
  and ordering comparisons at this chapter only accept
  \texttt{\textbf{int}}/\texttt{\textbf{float}} operands (see ``Dynamic
  Type Checking'' below), such an expression always raises a
  \texttt{\textbf{TypeError}} rather than silently returning an
  unintended value.
\item Unary \texttt{\textbf{+}}/\texttt{\textbf{-}}, \textbf{\texttt{not}}
  and \textbf{\texttt{lambda}} take a single operand to their right, so
  stacking two of them has only one possible grouping (e.g.\
  \texttt{\textbf{- -x}} always means \texttt{\textbf{-(-x)}});
  associativity does not apply and is marked ``n/a''.
\item The conditional expression and \textbf{\texttt{lambda}} have the
  lowest precedence and extend as far to the right as possible; a
  \textbf{\texttt{lambda}} body or the \textbf{\texttt{else}} branch of a
  conditional expression may itself be another conditional expression or
  \textbf{\texttt{lambda}}.
\end{itemize}
